\begin{table}[H]
\centering
\caption{Standardized Effect Sizes for Main Outcomes}
\label{tab:sde}
\small
\begin{tabular}{llccccc}
\toprule
Outcome & Spec. & $\hat{\beta}$ & SD($Y$) & SDE & SE(SDE) & Classification \\
\midrule
Employment & DDD & 0.0013 & 0.459 & 0.0027 & 0.0103 & Null \\
Hours/week & DDD & 0.011 & 18.59 & 0.0006 & 0.0108 & Null \\
LFP & DDD & 0.0047 & 0.435 & 0.0109 & 0.0117 & Small positive \\
\bottomrule
\end{tabular}
\par\vspace{0.3em}
\parbox{\textwidth}{\scriptsize
\textit{Notes:} \textbf{Country:} United States. \textbf{Research question:} Do state autism insurance mandates, which require private insurers to cover behavioral therapies for children with autism spectrum disorder, increase maternal labor force participation and employment? \textbf{Policy mechanism:} State-level mandates require private health insurers to cover diagnosis and treatment of autism spectrum disorder, including applied behavior analysis (ABA), speech therapy, and occupational therapy; by shifting the cost of intensive childhood therapy from families to insurers, mandates reduce the implicit caregiving tax on mothers who would otherwise provide or coordinate care themselves. \textbf{Outcome definition:} Employment is a binary indicator equal to one if the mother reports being employed (at work or with a job but not at work) in the ACS; hours per week is the usual weekly hours worked; labor force participation is a binary indicator for being employed or actively seeking work. \textbf{Treatment:} Binary---state adopted an autism insurance mandate (staggered across 46 states, 2001--2015). \textbf{Data:} American Community Survey 1-Year PUMS, 2008--2019, individual-level with household linkage; mothers (women aged 25--54, reference person or spouse) in households with children aged 5--17. \textbf{Method:} Triple-difference (state $\times$ year $\times$ child-disability-group) with state-group, year-group, and state-year fixed effects; standard errors clustered at the state level. \textbf{Sample:} Mothers aged 25--54 in households with school-age children; excludes states that adopted mandates before the ACS disability module began (2008). SDE $= \hat{\beta} / \text{SD}(Y)$ where SD($Y$) is the unconditional standard deviation from the full sample. Classification refers to magnitude, not statistical significance: Large ($|$SDE$| > 0.15$), Moderate ($0.05$--$0.15$), Small ($0.005$--$0.05$), Null ($< 0.005$).}
\end{table}
